%-------------------------
% Resume in LaTeX
% Author : Nilesh Rawal
%------------------------

\documentclass[letterpaper,11pt]{article}

\usepackage{latexsym}
\usepackage[empty]{fullpage}
\usepackage{titlesec}
\usepackage{marvosym}
\usepackage{enumitem}
\usepackage[hidelinks]{hyperref}
\usepackage{fancyhdr}

\pagestyle{fancy}
\fancyhf{}
\renewcommand{\headrulewidth}{0pt}
\renewcommand{\footrulewidth}{0pt}

% Adjust margins
\addtolength{\oddsidemargin}{-0.25in}
\addtolength{\textwidth}{0.5in}
\addtolength{\topmargin}{-.3in}
\addtolength{\textheight}{0.8in}

% Section formatting
\titleformat{\section}{
  \vspace{-4pt}\scshape\raggedright\large
}{}{0em}{}[\titlerule\vspace{-5pt}]

% Custom commands
\newcommand{\resumeItem}[1]{\item\small{#1}}
\newcommand{\resumeSubheading}[4]{
  \vspace{-2pt}\item
  \begin{tabular*}{0.97\textwidth}[t]{l@{\extracolsep{\fill}}r}
    \textbf{#1} & #2 \\
    \textit{\small#3} & \textit{\small #4} \\
  \end{tabular*}\vspace{-5pt}
}
\newcommand{\resumeSubHeadingListStart}{\begin{itemize}[leftmargin=0.15in, label={}]} 
\newcommand{\resumeSubHeadingListEnd}{\end{itemize}}
\newcommand{\resumeItemListStart}{\begin{itemize}[leftmargin=1.5em]} 
\newcommand{\resumeItemListEnd}{\end{itemize}\vspace{-5pt}}

%-------------------------------------------
%%%%%%  RESUME STARTS HERE  %%%%%%%%%%%%%%%%%%%%%%%%%%%%

\begin{document}

%----------HEADING-----------------
\begin{center}
    {\Huge \textbf{Nilesh Rawal}} \\[4pt]
    \href{mailto:rawalnilesh03@gmail.com}{rawalnilesh03@gmail.com} $|$
    +91 8237771166 $|$
    \href{https://linkedin.com/in/nileshrawal05}{linkedin.com/in/nileshrawal05} $|$
    \href{https://github.com/nileshrawal05}{github.com/nileshrawal05}
\end{center}

%-----------SUMMARY-----------------
\section{Summary}
Embedded Software Engineer with hands-on experience in platform bring-up, device driver development, and system integration across Embedded Linux, QNX, and Android on ARM64 and x86 platforms. Skilled in kernel-level debugging, network driver development (XGMAC on QNX), and Android display bring-up on virtual ARM SoCs. Experienced in diagnosing and resolving issues across the kernel-to-userspace stack, optimizing networking and system performance.
%-----------EXPERIENCE-----------------
\section{Experience}
\resumeSubHeadingListStart

\resumeSubheading
{Embedded Software Engineer}{Jul 2022 – Present}
{Vayavya Labs Pvt. Ltd.}{Belagavi, IN}

\resumeItemListStart
\resumeItem{Developed and integrated device drivers (Ethernet, UART, SPI) across QNX and Linux platforms on ARM64 and x86 architectures.}
\resumeItem{Contributed to Android platform bring-up on virtual ARM64 SoC, focusing on kernel-to-userspace debugging and issue resolution across JNI, DRM, and graphics layers.}
\resumeItem{Worked on EDMA (10G NIC) driver development across QNX and Linux, covering initialization, data path implementation, and performance optimization.}
\resumeItem{Build and customize Buildroot and Android build environments for platform adaptation.}
\resumeItem{Performed driver testing, debugging, and performance analysis using tools like GDB, strace, ftrace, and ktrace, ensuring stability and throughput optimization.}
\resumeItemListEnd

\resumeSubHeadingListEnd

%-----------KEY PROJECTS-----------------
\section{Key Projects}
\resumeItemListStart
\resumeItem{\textbf{Android Display \& Platform Bring-up – ARM Cortex Virtual SoC:}
Brought up Android display subsystem from kernel to HAL and userspace integration. Debugged kernel-level graphics pipeline and stabilized frame rendering and buffer management using \textbf{libdrm} and \textbf{OpenGL ES}. Integrated display stack with DRM, HAL, and Android framework layers. Customized Buildroot and AOSP builds to reproduce and debug boot and graphics failures. Performed low-level debugging (assembly, SoC timing, and configuration) to resolve hardware-level issues. Partnered with firmware and verification teams for validation and regression testing.}

\resumeItem{\textbf{EDMA / 10GbE Driver Development – QNX 7.1 (ARM64):}
Designed and implemented an EDMA-accelerated 10GbE Ethernet driver using \textbf{io-sock} and \textbf{iflib}. Ported core network stack mechanisms (interrupt handling via MSI-X, DMA descriptor management) from Linux to QNX. Implemented optimized buffer allocation and scatter-gather DMA for high throughput. Refactored driver for direct network stack integration, reducing software overhead. Validated DMA and interrupt performance using profiling and system-level testing on ARM64 hardware. Developed automated test suites to verify TX/RX data path efficiency and reliability.}
\resumeItemListEnd

%-----------TECHNICAL SKILLS-----------------
\section{Technical Skills}
\begin{itemize}[leftmargin=0.3in, label={}]
  \item \textbf{Languages:} C, C++, Embedded C, Shell Script
  \item \textbf{Platforms:} QNX, FreeBSD, Embedded Linux, Android (AOSP)
  \item \textbf{Drivers:} UART, Ethernet (io-sock, iflib), Display (libdrm, HAL)
  \item \textbf{Build Systems:} GCC, QCC, Buildroot, Yocto, U-Boot, AOSP
  \item \textbf{Tools:} Git, GDB, KGDB, strace, ftrace, perf, objdump
  \item \textbf{Debugging:} LLDB, Perfetto, ktrace, JTAG, Segger J-Link
  \item \textbf{Protocols:} UART, SPI, I2C, CAN, VLAN, ARP, TCP/IP
\end{itemize}

%-----------EDUCATION-----------------
\section{Education}
\resumeSubHeadingListStart
\resumeSubheading
{Post Graduate Diploma in Embedded Systems and Design}{2023}
{Centre for Development of Advanced Computing (C-DAC)}{Pune, IN}
\resumeSubheading
{Bachelor of Technology in Electrical Engineering}{2022}
{Sant Gajanan Maharaj College of Engineering}{Mahagaon, IN}
\resumeSubHeadingListEnd

\end{document}

